\documentclass[letterpaper, 11pt]{article}

\usepackage[utf8]{inputenc}
\usepackage[T1]{fontenc}
\usepackage[margin=1in, headheight=14pt]{geometry}
\usepackage{hyperref}
\usepackage{url}
\usepackage{booktabs}
\usepackage{array}
\usepackage{parskip}
\usepackage{xcolor}
\usepackage{fancyhdr}
\usepackage{listings}

\hypersetup{
  colorlinks=true,
  linkcolor=blue,
  urlcolor=blue,
}

\lstset{
  basicstyle=\ttfamily\small,
  breaklines=true,
  frame=single,
  backgroundcolor=\color{gray!8},
  xleftmargin=0.5em,
  xrightmargin=0.5em,
}

\pagestyle{fancy}
\fancyhf{}
\rhead{CS7641 --- Spring 2026}
\lhead{Reproducibility Sheet}
\rfoot{\thepage}

\title{\textbf{Reproducibility Sheet}\\[4pt]
\large CS7641: Machine Learning --- Spring 2026\\
Reinforcement Learning on Blackjack and CartPole}
\author{Siddartha Sagar Chinne \quad \texttt{schinne3@gatech.edu}}
\date{}

\begin{document}
\maketitle
\thispagestyle{fancy}

% ─────────────────────────────────────────────────────────────────────────────
\section*{1.\ Overleaf Link (READ-ONLY)}

\begin{center}
\url{https://www.overleaf.com/read/rnptvkjgxqpq#ece22b}
\end{center}

% ─────────────────────────────────────────────────────────────────────────────
\section*{2.\ GitHub Commit Hash}

The final code snapshot submitted with this report:

\begin{lstlisting}
0c84655818794185a99a9b319485f2d8f7f6cf4d
\end{lstlisting}

Repository hosted on GT Enterprise GitHub.

% ─────────────────────────────────────────────────────────────────────────────
\section*{3.\ Environment Setup}

\textbf{Python version:} 3.13 (managed via \href{https://docs.astral.sh/uv/}{uv})

\textbf{Operating system tested:} macOS 15 (Darwin); expected to run on standard Linux.

\textbf{No external datasets.} Both environments are provided by \texttt{gymnasium}.

\begin{lstlisting}
# Clone repository
git clone <gt-github-repo-url>
cd omscs_ml_rl

# Install uv (if not already installed)
curl -LsSf https://astral.sh/uv/install.sh | sh

# Create virtual environment and install all dependencies
make dev
\end{lstlisting}

\texttt{make dev} creates \texttt{.venv/} and installs all production and development
dependencies from \texttt{pyproject.toml} using \texttt{uv sync}.

% ─────────────────────────────────────────────────────────────────────────────
\section*{4.\ Reproduction Commands}

All phase scripts must be run via the \texttt{make} targets (which wrap
\texttt{bash ml\_run.sh} for sleep prevention). Do not invoke
\texttt{uv run python scripts/\ldots} directly for phase runs.

\subsection*{4.1\ Full Pipeline (recommended)}

\begin{lstlisting}
make pipeline
# Equivalent to: make phase1 gate1 phase2 phase3 phase4 phase5 phase6 phase8
\end{lstlisting}

\texttt{gate1} runs \texttt{pytest tests/test\_envs.py} after Phase 1 as a
correctness gate; Phase 2 will not start until it passes.

\subsection*{4.2\ Individual Phases}

\begin{lstlisting}
make phase1    # CartPole model estimation   ~40 s
make gate1     # Environment test gate       < 5 s  (required before Phase 2)
make phase2    # VI/PI on Blackjack          < 1 s
make phase3    # VI/PI on CartPole           ~70 s
make phase4    # SARSA/Q-Learning Blackjack  ~85 s
make phase5    # SARSA/Q-Learning CartPole   ~960 s
make phase6    # Cross-method comparison     < 5 s
make phase7    # DQN extra credit            ~95 s  (optional, not in pipeline)
make phase8    # Tables and repro artifacts  < 1 s
\end{lstlisting}

For phases expected to run longer than a few minutes, use the detached mode:

\begin{lstlisting}
bash ml_run.sh --detach "make phase5" phase5
\end{lstlisting}

\subsection*{4.3\ Figure Regeneration}

Figures are read from saved checkpoints and can be regenerated at any time
without rerunning experiments:

\begin{lstlisting}
make viz    # re-renders all phases
# or a single phase:
uv run python scripts/visualize_all.py phase5
\end{lstlisting}

% ─────────────────────────────────────────────────────────────────────────────
\section*{5.\ Random Seeds}

All model-free training runs (SARSA, Q-Learning, DQN) use the following five
seeds, defined as the constant \texttt{SEEDS} in \texttt{src/config.py}:

\begin{center}
\textbf{42, 43, 44, 45, 46}
\end{center}

Dynamic programming methods (VI, PI) are deterministic and require no seed.

% ─────────────────────────────────────────────────────────────────────────────
\section*{6.\ CartPole Discretization (Default Grid)}

The default discretization grid has \textbf{864 states} with bins
$(x, \dot{x}, \theta, \dot{\theta}) = (3, 3, 8, 12)$.
Non-uniform bin edges are listed verbatim below (from \texttt{src/config.py}).

\begin{center}
\renewcommand{\arraystretch}{1.15}
\begin{tabular}{lp{10cm}}
\toprule
\textbf{Dimension} & \textbf{Bin edges (inclusive)} \\
\midrule
$x$ (cart pos., m)
  & $[-2.4,\ -0.8,\ 0.8,\ 2.4]$ \\
$\dot{x}$ (cart vel., m/s)
  & $[-3.0,\ -1.0,\ 1.0,\ 3.0]$ \\
$\theta$ (pole angle, rad)
  & $[-0.200,\ -0.100,\ -0.050,\ -0.025,\ 0.000,\ 0.025,\ 0.050,\ 0.100,\ 0.200]$ \\
$\dot{\theta}$ (pole vel., rad/s)
  & $[-3.50,\ -2.00,\ -1.00,\ -0.50,\ -0.25,\ -0.10,\ 0.00,$\\
  & $\phantom{[}0.10,\ 0.25,\ 0.50,\ 1.00,\ 2.00,\ 3.50]$ \\
\bottomrule
\end{tabular}
\end{center}

Observations outside these ranges are clamped to the boundary bin.
Finer bins near zero on $\theta$ and $\dot{\theta}$ improve control resolution
where it matters most.

\medskip
Two additional grids are used in the discretization ablation study:

\begin{center}
\begin{tabular}{lccc}
\toprule
\textbf{Grid} & \textbf{Bins} & \textbf{States} & \textbf{Coverage} \\
\midrule
Coarse  & $(2, 2, 4, 6)$   & 96  & $\sim$41\% \\
Default & $(3, 3, 8, 12)$  & 864 & $\sim$83\% \\
Fine    & $(4, 4, 12, 16)$ & 3072 & $\sim$88\% \\
\bottomrule
\end{tabular}
\end{center}

% ─────────────────────────────────────────────────────────────────────────────
\section*{7.\ Key Hyperparameters}

\subsection*{Dynamic Programming}

\begin{tabular}{ll}
Discount factor $\gamma$ (VI and PI): & 0.99 \\
Convergence threshold $\delta$ (VI and PI): & $1 \times 10^{-6}$ \\
\end{tabular}

\subsection*{Model-Free (SARSA / Q-Learning)}

Hyperparameters were selected via a two-stage grid search.
Final tuned values (Blackjack and CartPole respectively) are recorded in
\texttt{artifacts/metadata/phase4.json} and \texttt{artifacts/metadata/phase5.json}.

\subsection*{DQN (Extra Credit)}

\begin{tabular}{ll}
Training episodes:          & 2{,}000 \\
Hidden layer dimension:     & 64 \\
Replay buffer size:         & 10{,}000 \\
Batch size:                 & 64 \\
Training start (steps):     & 1{,}000 \\
Target network update:      & every 500 steps \\
Learning rate:              & $5 \times 10^{-4}$ \\
$\varepsilon$ decay steps:  & 10{,}000 \\
$\gamma$:                   & 0.99 \\
\end{tabular}

% ─────────────────────────────────────────────────────────────────────────────
\section*{8.\ Expected Outputs}

All outputs are written under \texttt{artifacts/} (git-ignored):

\begin{center}
\begin{tabular}{ll}
\toprule
\textbf{Path} & \textbf{Description} \\
\midrule
\texttt{artifacts/metadata/phase\{N\}.json} & Checkpoint per phase \\
\texttt{artifacts/metrics/phase\{N\}\_*/}   & CSV metrics per phase \\
\texttt{artifacts/figures/phase\{N\}\_*/}   & PNG figures per phase \\
\texttt{artifacts/logs/phase\{N\}.log}      & Execution log per phase \\
\texttt{artifacts/tables/report\_numbers.tex} & Inline number macros \\
\texttt{artifacts/tables/tab\_*.tex}          & Table bodies \\
\texttt{artifacts/repro/runbook.md}           & Detailed runbook \\
\texttt{artifacts/repro/submission\_checklist.md} & Submission tracker \\
\bottomrule
\end{tabular}
\end{center}

\end{document}
